\section{信号怎样打断并恢复用户程序}

信号不是“内核替用户调用函数”。它先把异步事实记入 Process/Thread 状态，
再在安全的用户返回边界构造可验证的 Ring 3 控制流，并让条件、timeout 与
signal 竞争同一个 Blocked 到 Ready 提交点。

\topic{为什么 Unix 需要一种非轮询通知}
同步系统调用假定发起者知道自己在等待什么：读文件、等孩子或睡眠都由当前
执行流主动进入 Kernel。然而终端中断、非法指令、断开的管道、子进程退出和
管理员终止不会等待目标下一次查询。早期 Unix 为这些事件分配小整数编号，
并允许程序选择默认动作、忽略或用户处理函数，由此形成 signal。

signal 与硬件中断都能打破顺序控制流，但两者的特权边界不同。硬件 IRQ 由
CPU 根据 IDT 进入 CPL0，入口 frame 由 x86-64 架构定义；用户 signal 已经处于
Kernel 软件状态中，只在某个安全边界把即将恢复的 CPL3 RIP 改为 handler。
handler 从未在 PIC、PIT 或任意中断上下文中直接运行，也不会因“处理信号”
获得 Ring 0 权限。

\begin{pdfdiagram}[node distance=10mm and 13mm]
  \node[os hardware] (source) {发送者/异常/未来 TTY\\异步事件源};
  \node[os state,right=of source] (pending) {SignalManager\\pending 与处置};
  \node[os state,right=of pending] (thread) {选中一个 Thread\\或保留 Process pending};
  \node[os state,below=of thread] (boundary) {SYSCALL/IRQ\\用户返回边界};
  \node[os state,left=of boundary] (frame) {用户 SignalFrame\\原现场与 cookie};
  \node[os state,left=of frame] (handler) {Ring 3 handler\\restorer/sigreturn};
  \draw[os arrow] (source) -- (pending);
  \draw[os arrow] (pending) -- (thread);
  \draw[os arrow] (thread) -- (boundary);
  \draw[os arrow] (boundary) -- (frame);
  \draw[os arrow] (frame) -- (handler);
\end{pdfdiagram}
\diagramnote{异步事实先成为内核状态；真正改变用户 RIP 发生在统一返回边界。}

传统普通信号也不是消息队列。同一编号尚为 pending 时再次到达，只保证
“至少发生过一次”，并不积累次数。POSIX 后续增加实时信号队列以保存顺序和值；
本阶段实现普通合并语义，避免用一个无界计数伪装成尚未实现的实时信号。

\topic{处置、屏蔽和 pending 为什么属于不同所有者}
Process 是共享地址空间、文件和程序映像的资源容器；Thread 是独立寄存器、
栈和调度状态的执行实体。signal 横跨两层，不能把全部字段塞进一份“进程信号”
结构：

\begin{table}[htbp]
\centering
\small
\begin{osgridtabularx}{\textwidth}{l l X}
状态 & 所有者 & 不变量与设计原因\\ \hline
disposition & Process & 同一程序映像共享 Default/Ignore/Handler 决策\\ \hline
process group & Process & 广播目标身份，不等同于父子进程树\\ \hline
Process pending & Process & 尚无合格 Thread 时仍保存已经发生的事实\\ \hline
mask & Thread & 一个执行流的临界区不应屏蔽全部兄弟\\ \hline
Thread pending & Thread & 普通信号已经完成唯一接收者选择\\ \hline
active frame & Thread & sigreturn 只能恢复本 Thread 的寄存器与栈\\ \hline
\end{osgridtabularx}
\caption{信号状态分别保存在哪个对象中}
\end{table}

若 action 存在 Thread 中，新建 Thread 可能看到不同程序行为；若 mask 存在
Process 中，一个 Thread 暂时屏蔽就会意外推迟全部兄弟；若 active frame 只
存在 Process 中，兄弟 Thread 可能提交不属于自己的返回请求。
\code{SignalManager} 因而以固定容量 Process/Thread 表保存状态，再通过
scheduler 槽关联。公开 PID/TID 仍是单调身份，绝不能被数组槽号替代。

\topic{普通 pending 是位集合，不是次数}
当前有效信号号为 \(1\ldots63\)，编号 \(s\) 对应：

\[
  \operatorname{bit}(s)=1\ll(s-1)
\]

发送前同时查看目标 Process pending 与它全部 Thread pending。只要同一 bit
已经存在，本次发送增加 coalesced 统计但不产生第二个所有者。否则从 Process
保存的轮转游标扫描活动 Thread：

\begin{lstlisting}
thread.active
&& thread.process_index == target
&& (thread.mask & SignalBit(signal)) == 0
\end{lstlisting}

找到第一个合格者后只写它的 pending；没有合格者时写 Process pending。Thread
注册、修改 mask 或完成 handler 后重新分配 Process pending。轮转游标只消除
持续偏向最低槽位，不承诺实时优先级或公平期限。

\code{Kill} 是不可屏蔽、不可忽略、不可安装 handler 的信号。Kernel 在设置
mask 时强制清除对应 bit，并在 action 校验时拒绝非 Default；用户运行库的
检查只能改善错误报告，不能成为安全边界。

\topic{Default、Ignore 和 Handler 是三种提交}
\code{SignalAction} 固定为 40 字节，保存 disposition、handler、restorer、
additional mask 与 flags：

\begin{itemize}
  \item \textbf{Default}：Child 当前默认忽略；其他已定义信号终止 Process，
        退出原因记录为 Signal 并保存信号号；
  \item \textbf{Ignore}：消费 pending，不触碰用户栈，也不唤醒无关 Thread；
  \item \textbf{Handler}：构造用户 frame，把当前信号和 additional mask
        临时加入 Thread mask，再进入用户 handler。
\end{itemize}

handler 运行期间同一普通信号再次发送会保持 pending，因为本信号已自动
屏蔽。sigreturn 恢复旧 mask 后，下一次用户返回边界才可再投递。这阻止同一
信号无限递归堆叠；未来若支持 \code{nodefer}，必须新增明确的 action flag，
不能偷偷改变现有行为。

\topic{进程组不是进程树的另一种写法}
父子树回答谁 wait、谁收养 Zombie；进程组回答一次发送选择哪些 Process。
组是平坦成员集合，不包含“所有后代”的隐式递归。fork 继承父进程组，Process
可把自己移入以自身 PID 为组号的新组，按组发送再为每个活动成员独立执行一次
普通 Process signal。

整机 probe 必须先调用 \code{SetProcessGroup(0)} 建立独立组，再把 User2
设为 Ignore 并执行组发送。若 probe 留在 PID1 继承组，init 同样收到 User2
默认动作而终止；这不是随机投递故障，而是测试没有隔离广播身份。后续 TTY
将用前台进程组产生 Ctrl-C，但这里先只有身份和组发送，还没有 session、
控制终端或前后台作业。

\topic{signal 怎样加入单赢家唤醒}
Blocked Thread 已能同时关联业务 WaitQueue 与 deadline。信号路径不建立
绕过调度器的“信号就绪链”，而把 Signal 加入同一个完成竞争：

\begin{pdfdiagram}[node distance=12mm and 18mm]
  \node[os state] (thread) {Thread\\Blocked};
  \node[os state,below left=of thread] (condition) {condition/close\\普通通知};
  \node[os state,below=of thread] (timeout) {deadline\\Timeout};
  \node[os state,below right=of thread] (signal) {pending\\Signal};
  \node[os state,below=28mm of thread] (ready) {WakeThread\\Ready 一次};
  \draw[os arrow] (condition) -- (ready);
  \draw[os arrow] (timeout) -- (ready);
  \draw[os arrow] (signal) -- (ready);
  \draw[os arrow] (thread) -- node[right,font=\scriptsize]{同一 irq-save lock} (ready);
\end{pdfdiagram}
\diagramnote{三个原因都只能赢得一次 Blocked 到 Ready；失败者不覆盖 WakeReason。}

Signal 发送路径找到选中 Thread 后调用
\code{WakeThread(..., WakeReason::Signal)}。胜者在 scheduler irq-save lock
内确认仍为 Blocked、从 WaitQueue 摘除、取消 deadline、写入唯一 wake reason
并接入一次 Ready queue。condition 或 timeout 已先赢时，Signal 仍保持
pending，但不再次 Ready；Signal 先赢时，稍后的 deadline 已被计为 cancellation。

这条统一路径同时给出两个可审计不变量：

\[
  \text{一个阻塞周期只产生一个 WakeReason}
\]
\[
  active\ deadlines=schedules-expirations-cancellations
\]

\topic{Interrupted 与 restart 必须按调用语义区分}
信号唤醒一个阻塞系统调用时，Kernel 先把保存现场的 RAX 设置为固定错误
\code{-52}。是否在 handler 返回后重启，则同时取决于 action restart flag
和系统调用策略：

\begin{itemize}
  \item descriptor/pipe 读写、wait process 与 join 可重新检查条件；
  \item sleep 和带 deadline 的 futex 不重启；
  \item 已经提交部分字节的 I/O 必须返回部分结果，不能回退副作用。
\end{itemize}

对于允许重启且尚无部分提交的原生系统调用，signal frame 保存：

\[
  RIP_{\mathrm{saved}}=RIP_{\mathrm{after\ syscall}}-2
\]
\[
  RAX_{\mathrm{saved}}=\text{original syscall number}
\]

x86-64 \code{SYSCALL} 编码固定为两个字节；RDI、RSI、RDX、R10 等参数寄存器
仍在原 UserContext 中。sigreturn 恢复后，处理器重新执行同一条指令。sleep
若按相对时长盲目重启会把 handler 时间再次加入等待，futex 若重新比较错误
条件也会掩盖值变化，因此“可重启”不是一个全局布尔默认值。

\section{用户 handler 怎样进入和返回}
C++ handler 遵守 System V AMD64 ABI。普通函数入口期待
\(RSP\bmod16=8\)，栈顶保存返回地址；被中断代码却没有执行过 \code{call}。
Kernel 因而在目标用户栈合成：

\begin{lstlisting}
high address
  old RSP
  +-----------------------------+
  | SignalFrame, 240 bytes      | <- RSI
  +-----------------------------+  16-byte aligned
  | restorer return address     | <- handler RSP
low address
\end{lstlisting}

进入 handler 时 RDI 为信号号，RSI 为 frame 地址，RIP 为 action handler。
handler 执行 \code{ret} 后转入 Intel NASM restorer：

\begin{lstlisting}[language={[x86masm]Assembler}]
OsUserSignalReturnRestorer:
    mov rdi, rsp
    mov rax, 61
    syscall
    ud2
\end{lstlisting}

\code{ret} 已弹出 restorer 地址，此时 RSP 恰好指向 SignalFrame。成功
\code{SignalReturn} 直接恢复旧现场，不会普通返回；若 Kernel 意外返回，
\code{ud2} 把错误转成明确用户异常，避免沿未知栈继续执行。

\topic{SignalFrame 与 x86-64 标志位}
固定 240 字节 SignalFrame 包含 176 字节 UserContext、旧 mask、信号号、
magic/version/size、restorer 和 cookie。UserContext 中的 RFLAGS 不是任意
64 位用户数据。返回验证保留用户可用算术/控制状态，同时拒绝改变特权执行的位：

\begin{table}[htbp]
\centering
\small
\begin{osgridtabularx}{\textwidth}{l l X}
位/字段 & 处理 & 原因\\ \hline
CF/PF/AF/ZF/SF/OF & 恢复 & 普通整数运算结果属于用户现场\\ \hline
DF & 按统一白名单恢复 & 影响字符串指令方向，必须与通用返回契约一致\\ \hline
IF & 要求用户许可形态 & 不能让伪造 frame 破坏 Kernel 中断策略\\ \hline
IOPL/NT/VM & 拒绝危险值 & 不能通过 sigreturn 获得端口或旧任务/虚拟 8086 语义\\ \hline
RF/TF & 按项目白名单处理 & 调试与异常控制必须显式受控\\ \hline
CS/SS & 固定用户 selector & 返回只能进入 CPL3 代码/数据段\\ \hline
\end{osgridtabularx}
\caption{signal 返回复用统一 UserContext 的 RFLAGS 与段验证}
\end{table}

具体 bit mask 不在 signal 子系统另抄一份，而由统一
\code{ValidateUserContext} 契约判断。这样 SYSCALL 安全返回、异常返回和
sigreturn 不会各自演化出不同的 RFLAGS 漏洞。

\topic{主栈按需扩展与用户 Thread 栈边界}
主用户栈保留 8 MiB VMA，启动只提交参数真正占用的高端页面。SignalFrame 位于
当前 RSP 下方，可能跨入紧邻未提交页。
\code{PrepareUserStackRange} 验证 frame 与当前 RSP 均在主栈窗口，然后
只逐页提交紧邻 committed bottom 的页面，并复用用户写 not-present fault 的
VMA、权限和 64 KiB 邻近约束。

用户创建 Thread 的栈则是有 guard 的独立匿名 VMA。它不借用主栈增长几何；
Kernel 要求整个 frame 位于该 Thread 登记的 base/size 内，并确认页面用户可写。
fork 后若页面为 private COW，写入 frame 前仍经
\code{ResolveUserReturnMemory} 完成私有化。异步投递不能绕过 W\^X、guard
或 COW 所有权。

\topic{为什么 magic 之外还需要一次性 cookie}
用户能读写自己的栈，所以 magic、版本和大小只能发现随机损坏，不能证明这正是
Kernel 当前允许恢复的 frame。每次开始 handler 交付时，Kernel 在 Thread
状态登记精确五元组：

\begin{lstlisting}
frame_address
cookie
signal_number
restorer_address
previous_mask
\end{lstlisting}

frame 全部写入后才提交 active 登记。sigreturn 要求用户参数和五项逐一匹配，
成功后立即清空，所以旧 frame 不能重放，兄弟 Thread 也不能代为返回。cookie
是从非零开始的 64 位单调值，回绕遇零时跳过；它不是密码学秘密，但与精确
地址、Thread 所有权和一次消费共同构成能力凭据。

\topic{sigreturn 是特权返回验证器}
用户可以修改 frame 中的 RIP、RSP、RFLAGS 和段选择子。Kernel 若只
\code{CopyFromUser} 后直接 IRETQ，攻击者就能尝试返回 Kernel 高半地址、
不可执行页、错误段或危险 flags。返回路径依次验证：

\begin{enumerate}
  \item 当前 Thread 存在活动 frame；
  \item 参数地址与登记地址完全一致、16 字节对齐并完整位于本 Thread 栈；
  \item magic、version、size 与 reserved 字段正确；
  \item cookie、信号号、restorer 和 previous mask 与 Kernel 登记一致；
  \item 恢复 mask 不含不可屏蔽 Kill；
  \item RIP/RSP 满足四级页表 48 位 canonical 与用户低半范围；
  \item CS/SS 是固定用户 selector，RFLAGS 通过统一白名单；
  \item RIP 页为用户 R-X，RSP 页为用户 RW-NX。
\end{enumerate}

任何一步失败都复用非法用户返回隔离路径，只把目标 Process 记为
Exception/vector 13 并交给父进程 wait；Kernel 与旁系 Process 继续运行。
整机 probe 的孩子故意清零 magic，父进程必须观察隔离结果。这证明安全检查
穿过真实 ELF、页表和系统调用，而不是仅存在于宿主单元模型。

\section{fork、exec 与退出怎样迁移信号状态}
这三个生命周期动作都保留 Process 身份，却对信号状态采用不同的复制、
重置或销毁规则；把它们逐项展开，才能看清异步事件在哪一侧继续生效。

\topic{fork}
孩子复制 Process disposition 和进程组，唯一新 Thread 复制调用 Thread 的
mask；pending 与 active frame 不复制。否则父进程尚未处理的事件会在父子各
兑现一次，或孩子获得指向父用户栈的活动 frame。

\topic{exec}
新程序不知道旧 handler 地址，所以 Handler 重置为 Default，Ignore 保留。
Process pending 与 survivor pending 是已经发生的异步事实，和 survivor mask
一同保留；活动 frame 和被删除兄弟 Thread 的状态清空。早期实现曾错误清空
pending。TTY 整机测试曾让 Ctrl-Z 落入 fork/exec 窗口，由此抓到信号丢失；
由此修正。PID、父子关系与进程组没有变化，因为 exec 替换程序映像而非
Process 身份。

\topic{exit}
ThreadExit 移除对应信号状态；仍可交付的 Process pending 重新选择其他 Thread。
最后一个 Thread 或默认信号动作终止 Process 时，先终止 scheduler Thread，
再删除全部 ThreadSignalState，最后删除 ProcessSignalState。最终资源门禁要求
两类 active count 都为零。

\topic{固定 ABI、错误和值宽}
系统调用表在 0--56 号之后增加：

\begin{table}[htbp]
\centering
\small
\begin{osgridtabularx}{\textwidth}{l c X}
调用 & 编号 & 语义\\ \hline
\code{SetSignalAction} & 57 & 安装/读取 Process 处置\\ \hline
\code{SetSignalMask} & 58 & 替换当前 Thread mask 并返回旧值\\ \hline
\code{SendProcessSignal} & 59 & 向指定 PID 或自身发送普通信号\\ \hline
\code{SendProcessGroupSignal} & 60 & 向平坦进程组成员逐个发送\\ \hline
\code{SignalReturn} & 61 & 从精确活动 frame 恢复用户现场\\ \hline
\code{GetProcessGroup} & 62 & 读取 Process 组号\\ \hline
\code{SetProcessGroup} & 63 & 设置本 Process 组号\\ \hline
\end{osgridtabularx}
\caption{信号相关的固定宽度系统调用}
\end{table}

\code{Interrupted=-52}、\code{ProcessNotFound=-53} 与
\code{SignalStateInvalid=-54} 使用显式固定宽度枚举；action、context 与
frame 大小分别规定为 40、176 与 240 字节。ABI 不包含平台相关的
\code{long}、\code{size\_t} 或 Kernel 指针。

\topic{分别触发发送、递送、屏蔽和返回}
下面四层证据从局部状态机一路连接到真实用户程序；前一层负责扩大状态空间，
后一层负责确认真实硬件边界，任何一层都不能代替其余层。

\topic{单元测试}
直接驱动 SignalManager，验证两个 Thread mask 下唯一选择、普通信号合并、
active frame 五元组、fork 继承和 exec 重置。

\topic{集成测试}
组合 SignalManager、ThreadScheduler、WaitQueue 与 DeadlineQueue。Signal
先赢后，condition 再 wake、deadline 再 expire 都不能产生第二次 Ready，
deadline cancellation 统计必须守恒。

\topic{固定种子随机测试}
参考模型执行 100000 步发送、mask、投递和生命周期操作；每步比较 pending
并集，并要求同一 Process 的同一普通信号至多归属一个 Thread。种子与失败
步骤固定，可逐步重放。

\topic{QEMU 整机测试}
\filepath{/bin/signal\_probe} 穿过自研 ROM、Stage 1、rootfs ELF、原生
SYSCALL、用户页表、Intel NASM restorer、fork/COW 和 PID1 wait，依次证明：

\begin{itemize}
  \item descriptor wait 被 User1 打断并按 restart flag 重启；
  \item 屏蔽期间两次 User1 合并，解屏蔽后 handler 只执行一次；
  \item 独立进程组的 User2 组发送与 Ignore；
  \item fork 孩子继承 handler，COW 不污染父计数；
  \item 篡改 frame 只隔离提交它的孩子；
  \item Term 默认动作以 Signal 原因和信号号被父进程观察。
\end{itemize}

\topic{日志、故障定位与阶段边界}
发送和等待是热路径，Kernel 只在 action 安装、处置提交、失败以及累计计数
达到二次幂时输出局部日志，最终再打印 queued、coalesced、ignored、handler、
default termination、group send、interrupted/restarted wait 与 rejected
frame 汇总。用户 probe 只在一个语义事务完整提交后输出 marker；宿主为每行
添加 \code{[QEMU][T+...ms]} 到达时间，以识别静默卡死而不伪装成来宾时钟。

handler 未执行时先比较 action 和 delivery disposition，再检查 mask/pending
所有权与返回边界；handler 返回 vector 13 时检查 restorer 的 RSP、frame
cookie、用户 RIP 的 R-X 页和用户栈的 RW-NX 页；按组发送误伤 init 时先检查
是否仍继承 PID1 进程组。不能放宽 \code{ValidateUserContext} 或扩大 QEMU
timeout 来掩盖状态错误。

普通信号、进程组、用户 handler、严格 sigreturn 和可中断/可重启阻塞已经
具备。下面接入 TTY、session、前台组、Ctrl-C 与 stop/continue：i8042 输入
先经过 TTY line discipline，再由 TTY 向前台进程组发送信号；键盘 IRQ
不会直接选择某个 PID。alternate stack、实时信号队列和
\code{siginfo\_t} 尚未实现。
